\documentclass[aspectratio=169,11pt]{beamer}

\usetheme{Madrid}
\usecolortheme{default}
\setbeamertemplate{navigation symbols}{}

\usepackage[T1]{fontenc}
\usepackage[utf8]{inputenc}
\usepackage[english,ngerman]{babel}
\usepackage{amsmath,amssymb}
\usepackage{booktabs}
\usepackage{csquotes}
\usepackage{hyperref}
\usepackage{xcolor}
\usepackage{listings}

\lstdefinestyle{pycode}{
  language=Python,
  basicstyle=\ttfamily\scriptsize,
  keywordstyle=\color{blue!70!black},
  commentstyle=\color{green!40!black},
  stringstyle=\color{orange!60!black},
  showstringspaces=false,
  breaklines=true,
  frame=single,
  framerule=0.3pt,
  tabsize=2
}

\title{Use Evolutionary Strategies (ES) to fine-tune large language models (LLMs)}
\subtitle{Part I: Layer-Weight Optimization on GSM8K (my part)}
\author{Hasan Yildirim, Olsi Manco}
\institute{Uni Projekt PG ML/AI}
\date{19 Maerz, 2026}

\begin{document}

\begin{frame}
  \titlepage
\end{frame}

\begin{frame}{Presentation Structure}
\textbf{Short Introduction (shared, before Part I)}
\begin{itemize}
  \item Project context and motivation
  \item Why efficient LLM adaptation matters
\end{itemize}

\textbf{Part I (this talk segment): Hasan Yildirim}
\begin{itemize}
  \item ES-based fine-tuning via layer-weight perturbation
  \item Binary vs Log-Likelihood fitness and ES vs Backprop discussion
\end{itemize}

\textbf{Part II (following segment): Olsi Manco}
\begin{itemize}
  \item Soft prompting approach for LLM adaptation
  \item Comparison to layer-weight optimization
\end{itemize}
\end{frame}

\begin{frame}{Short Introduction (before Part I)}
\textbf{Project context}
\begin{itemize}
  \item This presentation studies two parameter-efficient adaptation paths for LLMs.
  \item Part I focuses on Evolutionary Strategies (ES) over layer weights.
  \item Part II focuses on soft prompting as an alternative adaptation mechanism.
\end{itemize}

\textbf{Bridge to my part}
\begin{itemize}
  \item Key question for Part I: can we improve reasoning performance using a gradient-free optimizer on selected layers?
\end{itemize}
\end{frame}

\begin{frame}{Primer 1: What is Fine-Tuning?}
\textbf{Simple definition}
\begin{itemize}
  \item A pre-trained LLM already has broad language knowledge.
  \item Fine-tuning means adapting this model to a specific task or domain.
\end{itemize}

\textbf{In this project}
\begin{itemize}
  \item Task: mathematical reasoning (GSM8K).
  \item We do not retrain the full model from scratch.
  \item We adapt selected model parameters to improve task performance.
\end{itemize}
\end{frame}

\begin{frame}{Primer 2: Why not only Backpropagation?}
\textbf{Backpropagation is strong, but not always ideal for our objective space}
\begin{itemize}
  \item Backprop needs differentiable training signals.
  \item Some practical goals are discrete/non-smooth (e.g., exact-match 0/1).
  \item We also want to test a black-box optimizer that only needs a scalar score.
\end{itemize}

\textbf{Research angle}
\begin{itemize}
  \item Can we still improve LLM behavior when we optimize without explicit gradients?
\end{itemize}
\end{frame}

\begin{frame}{Primer 3: ES as a Gradient-Free Alternative}
\textbf{High-level idea of Evolutionary Strategies (ES)}
\begin{enumerate}
  \item Start from current parameter vector $\theta$.
  \item Create several noisy variants ($\theta + \epsilon_i$).
  \item Evaluate each variant with a fitness score.
  \item Move $\theta$ toward better-performing variants.
\end{enumerate}

\textbf{Key intuition}
\begin{itemize}
  \item Backprop asks: \enquote{What is the exact gradient?}
  \item ES asks: \enquote{Which parameter perturbations worked better?}
\end{itemize}
\end{frame}

\begin{frame}{Problem Setting: Backpropagation vs ES}
\textbf{Core problem in this project}
\begin{itemize}
  \item We want to optimize task quality on GSM8K, where a natural metric is exact-match (0/1).
  \item Exact-match is discrete and not smoothly differentiable.
  \item Standard backpropagation is strongest when we optimize differentiable objectives directly.
\end{itemize}

\textbf{Why compare with ES?}
\begin{itemize}
  \item ES only needs a scalar fitness value per candidate.
  \item ES can therefore optimize with non-differentiable or noisy task-level signals.
  \item This gives a practical alternative when gradient quality is weak or objective mismatch is high.
\end{itemize}
\end{frame}

\begin{frame}{Challenges of the Problem}
\textbf{Optimization challenges}
\begin{itemize}
  \item High-dimensional parameter space, even with layer-wise selection.
  \item Noisy fitness estimates due to sampling and finite evaluation subsets.
  \item Trade-off between fast surrogate fitness (LogLik) and final target metric (Exact-Match).
\end{itemize}

\textbf{Systems and evaluation challenges}
\begin{itemize}
  \item ES is evaluation-heavy: population size $\times$ samples per generation.
  \item Binary evaluation is expensive because of autoregressive generation.
  \item Fair comparison needs strict protocol control (same seeds, same limits, same splits).
\end{itemize}
\end{frame}

\begin{frame}{Motivation}
\begin{itemize}
  \item Klassisches LLM-Finetuning basiert auf Backpropagation und dichten Gradienten.
  \item Viele praxisnahe Zielgroessen sind jedoch diskret oder nicht glatt (z.\,B. Exact-Match auf Aufgabenebene).
  \item Forschungsfrage: Wie weit kommen wir mit \textbf{gradient-freiem} Finetuning via Evolutionary Strategies (ES)?
\end{itemize}
\vspace{0.5em}
\textbf{Use Case:} Mathematisches Reasoning (GSM8K) mit einem kleinen Instruct-LLM.
\end{frame}

\begin{frame}{Projektziel und Forschungsfragen}
\textbf{Ziel:} Verbesserung der Task-Performance durch selektive Anpassung weniger Modellparameter.

\vspace{0.6em}
\textbf{Forschungsfragen}
\begin{enumerate}
  \item Kann ES auf wenigen Decoder-Layern sinnvolle Verbesserungen liefern?
  \item Welche Fitness-Definition ist fuer ES stabil und effizient?
  \item Wie verhaelt sich ES konzeptionell zu Backpropagation in diesem Setting?
\end{enumerate}
\end{frame}

\begin{frame}{Modellvorstellung und Datensetup}
\textbf{Basismodell}
\begin{itemize}
  \item \texttt{Qwen/Qwen2.5-0.5B-Instruct}
  \item Causal LM, 24 Decoder-Layer, Instruct-Chat-Format
\end{itemize}

\textbf{Task}
\begin{itemize}
  \item GSM8K (arithmetisches Reasoning)
  \item Training-Fitness auf \texttt{train}-Split, finale Bewertung auf \texttt{test}-Split
\end{itemize}

\textbf{Projektarchitektur}
\begin{itemize}
  \item Layer-Selektion: \texttt{src/es\_llm/model/layer\_selector.py}
  \item ES-Optimierer: \texttt{src/es\_llm/es/openai\_es.py}
  \item Fitness-Module: \texttt{gsm8k.py}, \texttt{gsm8k\_loglikelihood.py}
\end{itemize}
\end{frame}

\begin{frame}{Task Primer: GSM8K vs HellaSwag (mit Beispiel)}
\textbf{GSM8K (Grade School Math 8K)}
\begin{itemize}
  \item Ziel: mehrschrittige Rechenaufgaben loesen und finale Zahl liefern.
  \item Beispiel: \enquote{Sara has 12 apples, buys 5, gives away 4. How many remain?}
  \item Erwartung: Schrittweises Reasoning, finale Antwort z.\,B. \texttt{#### 13}.
\end{itemize}

\textbf{HellaSwag}
\begin{itemize}
  \item Ziel: die plausibelste Fortsetzung einer Situation aus mehreren Optionen waehlen.
  \item Beispiel: Kontext \enquote{A person cracks eggs into a bowl...}
  \item Optionen: (A) plausible Koch-Fortsetzung, (B/C/D) unplausible Ablenkungen.
\end{itemize}

\textbf{Unterschied kurz:}
\begin{itemize}
  \item GSM8K misst numerisches/mathematisches Reasoning.
  \item HellaSwag misst commonsense narrative completion.
\end{itemize}
\end{frame}

\begin{frame}{Qwen2.5-0.5B: Gesamtstruktur (visuell)}
\centering
\small
\textbf{Token Input}

\vspace{0.25em}
$\downarrow$

\textbf{Embedding Layer: \texttt{model.embed\_tokens}}

\vspace{0.25em}
$\downarrow$

\textbf{Decoder Stack: Layer 1, 2, \dots, 24}

\vspace{0.25em}
$\downarrow$

\textbf{Final Norm: \texttt{model.norm}}

\vspace{0.25em}
$\downarrow$

\textbf{Output Head: \texttt{lm\_head}}

\vspace{0.8em}
\begin{columns}[T]
\column{0.5\textwidth}
\textbf{Layerbereiche}
\begin{itemize}
  \item Fruehe Layer: 1--8
  \item Mittlere Layer: 9--16
  \item Spaete Layer: 17--24
\end{itemize}

\column{0.5\textwidth}
\textbf{Parameterfamilien im Modell}
\begin{itemize}
  \item Embedding-Parameter
  \item Self-Attention-Parameter
  \item MLP-Parameter
  \item LayerNorm-Parameter
  \item Output-Head-Parameter
\end{itemize}
\end{columns}
\end{frame}

\begin{frame}{Aufbau eines einzelnen Decoder-Layers}
\textbf{Pro Layer (z. B. Layer 20) werden strukturell diese Bl"ocke genutzt:}

\vspace{0.4em}
\begin{columns}[T]
\column{0.5\textwidth}
\textbf{1) Attention-Block}
\begin{itemize}
  \item \texttt{self\_attn.q\_proj.weight}
  \item \texttt{self\_attn.k\_proj.weight}
  \item \texttt{self\_attn.v\_proj.weight}
  \item \texttt{self\_attn.o\_proj.weight}
\end{itemize}

\textbf{2) Norms}
\begin{itemize}
  \item \texttt{input\_layernorm.weight}
  \item \texttt{post\_attention\_layernorm.weight}
\end{itemize}

\column{0.5\textwidth}
\textbf{3) MLP-Block}
\begin{itemize}
  \item \texttt{mlp.gate\_proj.weight}
  \item \texttt{mlp.up\_proj.weight}
  \item \texttt{mlp.down\_proj.weight}
\end{itemize}

\textbf{Fuer dieses Projekt wichtig}
\begin{itemize}
  \item Layer-wise Auswahl ueber \texttt{layer\_indices}
  \item Komponentenwahl ueber \texttt{components}
  \item Beispiel: nur \texttt{attention} in Layer 18--21
\end{itemize}
\end{columns}
\end{frame}

\begin{frame}{Warum Layer-wise statt Full Finetuning?}
\begin{itemize}
  \item Reduziert Suchraum und Compute-Kosten.
  \item Ermoeglicht kontrollierte Experimente zu \enquote{welche Layer tragen am meisten bei?}
  \item Passt zu ES: kleinere, zielgerichtete Parametervektoren sind robuster optimierbar.
\end{itemize}

\textbf{Konkretes Setup im Projekt}
\begin{itemize}
  \item Auswahl nach Layer-Indizes (z.\,B. \texttt{[18,19,20,21]}).
  \item Komponentenfokus, z.\,B. nur \texttt{attention} statt \texttt{all}.
\end{itemize}
\end{frame}

\begin{frame}{ES-Algorithmus (OpenAI-ES)}
\textbf{Ablauf pro Generation}
\begin{enumerate}
  \item \textbf{ask:} Stichprobe von Perturbationen $\epsilon_i \sim \mathcal{N}(0,\sigma^2 I)$
  \item Kandidaten evaluieren: $F(\theta + \epsilon_i)$
  \item \textbf{tell:} Parameterupdate auf Basis eines Gradienten-Schaetzers
\end{enumerate}

\vspace{0.4em}
\[
\hat{g} = \frac{1}{N\sigma}\sum_{i=1}^N u_i\,\epsilon_i,
\qquad
\theta_{t+1}=\theta_t + \alpha\hat{g}
\]

\textbf{Im Projekt aktiviert:}
\begin{itemize}
  \item Antithetic Sampling
  \item Fitness shaping (centered rank)
  \item Optional Adam-Style Update
  \item Sigma-Decay (z.\,B. cosine)
\end{itemize}
\end{frame}

\begin{frame}[fragile]{Pseudocode: Evolution Strategies (Salimans et al., Alg. 1)}
\small
\begin{lstlisting}[style=pycode]
Algorithm: Evolution Strategies

Input:
  learning rate alpha,
  noise standard deviation sigma,
  initial parameters theta_0

for t = 0, 1, 2, ... do
  Sample eps_1, ..., eps_n ~ N(0, I)

  Compute returns:
    F_i = F(theta_t + sigma * eps_i),  i = 1, ..., n

  Update parameters:
    theta_{t+1} <- theta_t + alpha * (1 / (n * sigma)) *
                   sum_{i=1..n} (F_i * eps_i)
end for
\end{lstlisting}

\textbf{Interpretation:} Der Update-Schritt ist ein stochastischer Gradienten-Schaetzer aus Funktionsauswertungen.
\end{frame}

\begin{frame}[fragile]{Code: ES \texttt{ask()} und Kandidatenbildung}
\textbf{Wichtig fuer die Methodik:} Kandidaten entstehen aus Perturbationen, nicht aus Gradienten.

\begin{lstlisting}[style=pycode]
# src/es_llm/es/openai_es.py
def ask(self, current_params):
  sigma = self._get_sigma()
  if self.antithetic:
    noises = antithetic_pairs(dim, n_pairs, sigma=sigma)
  else:
    noises = [gaussian_noise(dim, sigma=sigma)
          for _ in range(self.population_size)]

  self._noise_cache = noises
  candidates = [current_params + eps for eps in noises]
  return candidates
\end{lstlisting}

\textbf{Takeaway:} ES sampelt den Suchraum explizit und evaluiert Black-Box-Fitness auf jedem Kandidaten.
\end{frame}

\begin{frame}[fragile]{Code: ES \texttt{tell()} Update ohne \texttt{backward()}}
\textbf{Wichtig fuer die Argumentation ES vs Backprop:}

\begin{lstlisting}[style=pycode]
# src/es_llm/es/openai_es.py
def tell(self, current_params, candidates, fitnesses):
  shaped = _centered_rank_transform(fitnesses)

  grad = torch.zeros_like(current_params, dtype=torch.float32)
  for i, eps in enumerate(self._noise_cache):
    grad += float(shaped[i]) * eps.float()
  grad /= len(self._noise_cache) * self._cached_sigma

  new_params = current_params.float() + self.lr * grad
  return ESResult(new_params=new_params.to(current_params.dtype), ...)
\end{lstlisting}

\textbf{Takeaway:} Der naechste Parametervektor kommt aus Fitness-gewichteten Stoerungen, nicht aus Autograd-Gradienten.
\end{frame}

\begin{frame}{Backpropagation vs Evolutionary Strategies}
\begin{columns}[T]
\column{0.49\textwidth}
\textbf{Backpropagation}
\begin{itemize}
  \item Nutzt exakte Ableitungen $\nabla_\theta L$
  \item Effizient bei differenzierbaren Losses
  \item Braucht Autograd-Graph und \texttt{backward()}
\end{itemize}

\column{0.49\textwidth}
\textbf{Evolutionary Strategies}
\begin{itemize}
  \item Braucht nur Funktionswerte $F(\theta)$
  \item Funktioniert auch fuer nicht-differenzierbare Ziele
  \item Kein \texttt{backward()}, kein expliziter Gradient notwendig
\end{itemize}
\end{columns}

\vspace{0.6em}
\textbf{Kernaussage:} ES ist ein Black-Box-Optimierer; Backprop ist ein Gradient-Optimierer.
\end{frame}

\begin{frame}{Anlauf 1: Binary Validator als Fitness}
\textbf{Idee}
\begin{itemize}
  \item Prompt $\rightarrow$ \texttt{model.generate(...)}
  \item Antwort extrahieren (\texttt{#### <answer>})
  \item Fitness = Exact-Match (0/1)
\end{itemize}

\textbf{Staerken}
\begin{itemize}
  \item Zielmetrik-nah (entspricht Endbewertung)
  \item Klar interpretierbar
\end{itemize}

\textbf{Schwaechen}
\begin{itemize}
  \item Nicht-differenzierbar und diskret
  \item Hohe Varianz im ES-Signal
  \item Langsam durch autoregressive Generierung
\end{itemize}
\end{frame}

\begin{frame}{Anlauf 2: Log-Likelihood als Fitness}
\textbf{Idee}
\begin{itemize}
  \item Keine freie Generierung, sondern Scoring der Gold-Token
  \item Fitness = mittlere Log-Wahrscheinlichkeit der korrekten Antworttokens
\end{itemize}

\textbf{Warum schneller?}
\begin{itemize}
  \item Ein Forward-Pass (bzw. batched Forward-Paesse) statt Token-fuer-Token-Decoding
  \item Caching der tokenisierten Inputs
\end{itemize}

\textbf{Warum stabiler fuer ES?}
\begin{itemize}
  \item Kontinuierliches Signal statt 0/1-Sprung
  \item Feineres Ranking zwischen Kandidaten in einer Population
\end{itemize}
\end{frame}

\begin{frame}{Log-Likelihood Algorithmus (Schritt fuer Schritt)}
\textbf{Ziel:} Bewerte, wie wahrscheinlich das Modell die \emph{korrekte} Antwort findet.

\vspace{0.4em}
\textbf{Ablauf pro Beispiel}
\begin{enumerate}
  \item Erzeuge Prompt aus \texttt{system} + \texttt{user} Nachricht.
  \item Haenge Gold-Antwort an: \texttt{full\_text = prompt + gold\_answer}.
  \item Tokenisiere und merke \texttt{answer\_start} (Start der Antwort-Tokens).
  \item Forward-Pass: \texttt{logits = model(input\_ids)}.
  \item Wandle in Log-Wahrscheinlichkeiten um: \texttt{log\_softmax(logits)}.
  \item Sammle nur die Gold-Token-LogProbs im Antwortbereich.
  \item Mittelwert ueber Antwort-Tokens und dann ueber alle Samples.
\end{enumerate}

\vspace{0.4em}
\textbf{Formel (vereinfacht)}
\[
\text{fitness} = \frac{1}{M} \sum_{j=1}^{M}
\left(
\frac{1}{T_j} \sum_{t=1}^{T_j}
\log p_\theta\big(y_{j,t}^{\text{gold}} \mid x_j, y_{j,<t}^{\text{gold}}\big)
\right)
\]

\textbf{Interpretation:} Hoeher (weniger negativ) bedeutet: Modell weist den korrekten Tokens mehr Wahrscheinlichkeit zu.
\end{frame}

\begin{frame}[fragile]{Code: Binary vs LogLik Fitnesspfad}
\textbf{Binary (langsam, diskret):}
\begin{lstlisting}[style=pycode]
# src/es_llm/fitness/gsm8k.py
gen = model.generate(**inputs, max_new_tokens=...)
pred_text = tokenizer.decode(out_ids, skip_special_tokens=True)
pred = extract_hash_answer(pred_text)
fitness = int(pred == gold)   # 0/1
\end{lstlisting}

\textbf{LogLik (schnell, kontinuierlich):}
\begin{lstlisting}[style=pycode]
# src/es_llm/fitness/gsm8k_loglikelihood.py
logits = model(input_ids, attention_mask=attention_mask).logits
log_probs = torch.nn.functional.log_softmax(logits, dim=-1)
token_lls = ans_log_probs.gather(1, ans_tokens.unsqueeze(1)).squeeze(1)
fitness = token_lls.mean().item()
\end{lstlisting}

\textbf{Takeaway:} Beide liefern nur einen Skalar an ES; der Rechenpfad bestimmt Laufzeit und Signalqualitaet.
\end{frame}

\begin{frame}[fragile]{Code: Trainingsloop (Orchestrierung)}
\textbf{Projektkern in \texttt{es\_trainer.py}:}

\begin{lstlisting}[style=pycode]
for gen in range(1, num_generations + 1):
  fitness.reshuffle_data()
  candidates = es.ask(current_params)

  fitnesses = []
  for cand in candidates:
    selector.set_flat_params(cand)
    fit = fitness.evaluate(model, tokenizer)
    fitnesses.append(fit)

  result = es.tell(current_params, candidates, fitnesses)
  current_params = result.new_params.clone()
\end{lstlisting}

\textbf{Takeaway:} \texttt{ask -> evaluate -> tell} ist die wiederholte Optimierungsschleife ueber alle Generationen.
\end{frame}

\begin{frame}{Wichtige konzeptionelle Klarstellung}
\textbf{Ist Log-Likelihood differenzierbar?} Ja.

\vspace{0.4em}
\textbf{Nutzen wir diesen Gradient?} Nein, in diesem Projekt nicht.
\begin{itemize}
  \item Fitness wird in \texttt{inference\_mode()} berechnet.
  \item Kein \texttt{loss.backward()} im Trainingsloop.
  \item Naechster Parametervektor kommt aus ES-\texttt{ask/tell}, nicht aus Autograd.
\end{itemize}

\vspace{0.5em}
\textbf{Interpretation}
\begin{itemize}
  \item Binary zeigt, dass ES auf nicht-differenzierbaren Zielen funktioniert.
  \item LogLik zeigt, dass ES auch auf differenzierbaren Surrogates praktikabel und schneller sein kann.
\end{itemize}
\end{frame}

\begin{frame}{Experimentelles Setup (Beispielkonfiguration)}
\textbf{Konfiguration: \texttt{gsm8k\_a100\_v5.yaml}}
\begin{itemize}
  \item Layer: \texttt{[18,19,20,21]}, Komponenten: \texttt{attention}
  \item ES: Population 96, 100 Generationen, $\sigma:0.003 \rightarrow 0.001$ (cosine)
  \item Fitness: \texttt{gsm8k\_loglikelihood}, \texttt{num\_samples=120}, \texttt{batch\_size=12}
\end{itemize}

\textbf{Pipeline}
\begin{itemize}
  \item Baseline-Fitness
  \item Pro Generation: \texttt{ask -> evaluate(population) -> tell}
  \item Checkpointing und Logging pro Run
\end{itemize}
\end{frame}

\begin{frame}{Bisherige Beobachtungen und Learnings}
\begin{itemize}
  \item Binary-Fitness ist methodisch nah an Endmetrik, aber teuer und noisig.
  \item LogLik-Fitness ist deutlich effizienter und liefert glattere Optimierungssignale.
  \item Fitnessdefinition dominiert die reale Laufzeit im ES-Loop
  \item Layer-selektive Optimierung ist praktikabel fuer kontrollierte Studien.
\end{itemize}

\vspace{0.4em}
\textbf{Wissenschaftliche Konsequenz:}
\begin{itemize}
  \item Fuer faire Aussagekraft sind strikte A/B-Setups noetig (gleicher Seed, gleiches Eval-Budget).
\end{itemize}
\end{frame}

\begin{frame}{Threats to Validity}
\begin{itemize}
  \item Unterschiedliche Eval-Limits/Splits koennen Effekte verzerren.
  \item Surrogate mismatch: Trainingsfitness (LogLik) vs Zielmetrik (Exact-Match).
  \item Seed-Sensitivitaet bei stochastischer Populationsevaluation.
  \item Moegliches Overfitting auf Teilmengen ohne sauberes Reshuffling-Protokoll.
\end{itemize}
\end{frame}

\begin{frame}{Fazit und Ausblick}
\textbf{Fazit}
\begin{itemize}
  \item ES ist ein valider Ansatz fuer gradient-freies LLM-Finetuning auf selektierten Layern.
  \item Besonders starkes Argument bei nicht-differenzierbaren Zielgroessen (Binary Validator).
  \item LogLik als Fitness ist ein effizienter, stabiler Surrogate fuer ES-Training.
\end{itemize}

\textbf{Ausblick}
\begin{itemize}
  \item Saubere Baseline-vs-ES Studien mit mehreren Seeds
  \item Ablationen: Layerauswahl, Population, Sigma-Schedule
  \item Transfer-Evaluierung auf weitere Reasoning-Tasks
\end{itemize}
\end{frame}

\begin{frame}{Backup: Reproduzierbarkeit}
\textbf{Run-Artefakte}
\begin{itemize}
  \item \texttt{config.json}
  \item \texttt{training\_log.jsonl}
  \item \texttt{checkpoint\_genXXXX.pt}
  \item \texttt{best\_layer\_weights.pt}
  \item gespeichertes Modell unter \texttt{model/}
\end{itemize}

\textbf{Entry Points}
\begin{itemize}
  \item Training: \texttt{scripts/train\_es.py}
  \item Evaluierung: \texttt{scripts/eval\_es\_model.py}, \texttt{scripts/eval\_olmes\_tasks.py}
\end{itemize}
\end{frame}

\end{document}
